\chapter{Boolean Analysis — Centered}
%%% generated by the blueprint pipeline from TCSlib.BooleanAnalysis.Hypercontractivity.Randomization.Centered
\section{Overview}
Absolute-power remainders yield a uniform scalar reflected contraction for every exponent at least two. Its affine form gives norm contraction for centered random variables on probability spaces.

\section{Declarations}

\begin{theorem}[Quadratic remainder lower bound]
\lean{BooleanAnalysis.Hypercontractivity.CenteredContraction.remainder_lower}
\leanok
\difficulty{4}
For every $q\ge2$ and every $x\in\mathbb R$,
\[
x^2\le |1+x|^q-1-qx.
\]
\end{theorem}

\begin{theorem}[Integer-power remainder bound]
\lean{BooleanAnalysis.Hypercontractivity.CenteredContraction.pow_remainder_bound}
\leanok
\difficulty{4}
For every natural number $m$ and every real $x$ with $|x|\le1$,
\[
(1+x)^m-1-mx\le (2^m-1-m)x^2.
\]
\end{theorem}

\begin{theorem}[Local real-power remainder bound]
\lean{BooleanAnalysis.Hypercontractivity.CenteredContraction.remainder_local_upper}
\leanok
\uses{BooleanAnalysis.Hypercontractivity.CenteredContraction.pow_remainder_bound}
\difficulty{5}
Let $q\ge2$, let $m\in\mathbb N$ satisfy $q\le m$, and let $x\in\mathbb R$ satisfy
$|x|\le1$. Then
\[
|1+x|^q-1-qx\le 2^m x^2.
\]
\end{theorem}

\begin{theorem}[Global remainder bound]
\lean{BooleanAnalysis.Hypercontractivity.CenteredContraction.remainder_upper}
\leanok
\uses{BooleanAnalysis.Hypercontractivity.CenteredContraction.remainder_local_upper}
\difficulty{5}
For every real $q\ge 2$, there is a constant $M>0$ such that, for every real $x$,
$|1+x|^q-1-qx\le M(x^2+|x|^q)$.
\end{theorem}

\begin{theorem}[Translated absolute-power bound]
\lean{BooleanAnalysis.Hypercontractivity.CenteredContraction.translated_power_bound}
\leanok
\difficulty{4}
For every real $q\ge 0$ and every real $x$, $|x|^q\le 2^q(|1+x|^q+1)$.
\end{theorem}

\begin{theorem}[Power controlled by affine remainder]
\lean{BooleanAnalysis.Hypercontractivity.CenteredContraction.remainder_power_lower}
\leanok
\uses{BooleanAnalysis.Hypercontractivity.CenteredContraction.remainder_lower, BooleanAnalysis.Hypercontractivity.CenteredContraction.translated_power_bound}
\difficulty{5}
For every $q\ge2$ and every $x\in\mathbb R$,
\[
|x|^q\le 2^q(q+3)\bigl(|1+x|^q-1-qx\bigr).
\]
\end{theorem}

\begin{theorem}[Uniform scalar reflected contraction]
\lean{BooleanAnalysis.Hypercontractivity.CenteredContraction.scalar_contraction_generic}
\leanok
\uses{BooleanAnalysis.Hypercontractivity.CenteredContraction.remainder_lower, BooleanAnalysis.Hypercontractivity.CenteredContraction.remainder_power_lower, BooleanAnalysis.Hypercontractivity.CenteredContraction.remainder_upper}
\difficulty{5}
For every real $q\ge 2$, there is a real $c$ with $0<c\le 1$ such that, for every real
$x$, $|1-cx|^q\le |1+x|^q-q(1+c)x$.
\end{theorem}

\begin{theorem}[Reflected scalar contraction]
\lean{BooleanAnalysis.Hypercontractivity.CenteredContraction.scalar_contraction}
\leanok
\uses{BooleanAnalysis.Hypercontractivity.CenteredContraction.scalar_contraction_generic}
\difficulty{4}
For every real $q\ge 2$, there is a real $c$ with $0<c\le 1$ such that, for every real
$x$, $|1-cx|^q\le |1+x|^q-q(1+c)x$. The constant can be chosen to satisfy $c=2/5$ when
$q=4$.
\end{theorem}

\begin{theorem}[Affine scalar contraction]
\lean{BooleanAnalysis.Hypercontractivity.CenteredContraction.scalar_affine}
\leanok
\difficulty{4}
Let $q>0$ and $0\le c\le 1$, and suppose that $|1-cz|^q\le |1+z|^q-q(1+c)z$ for every
real $z$. Then, for all real $a$ and $x$, $|a-cx|^q\le |a+x|^q-[q(1+c)|a|^q/a]x$, where
the quotient is interpreted as zero when $a=0$.
\end{theorem}

\begin{theorem}[Centered random-variable norm contraction]
\lean{BooleanAnalysis.Hypercontractivity.CenteredContraction.norm_le_of_scalar}
\leanok
\uses{BooleanAnalysis.Hypercontractivity.CenteredContraction.scalar_affine, BooleanAnalysis.Hypercontractivity.rvLpNorm, BooleanAnalysis.Hypercontractivity.rvLpNorm_rpow}
\difficulty{5}
Let $(\Omega,\mu)$ be a probability space, let $q\ge2$, and let $c\in[0,1]$. Suppose
that for every $z\in\mathbb R$,
\[
|1-cz|^q\le |1+z|^q-q(1+c)z.
\]
If $X\in L^q(\mu)$ has mean zero, then for every $a\in\mathbb R$, the real-valued
random-variable norms defined by converting the extended $L^q$ norm to a real number
satisfy
\[
N_q(a-cX)\le N_q(a+X).
\]
\end{theorem}
